% Source PDF: source/普通高考/2009/2009安徽理.pdf, page 1, question 13.
% Grid evidence: tmp/review_flowchart_2009_anhui_science/grid/page-1-grid100.png and
%   q13-block-grid50.png; final comparison crop is from the no-grid page render.
% Node-edge list: 开始→a=1→a=2a+1→a>100?; 是→输出a→结束;
% 否→左回边→a=2a+1. The loop return joins the main stem above a=2a+1.
% No other edges, labels, or dashed connectors occur in the source figure.

\begin{tikzpicture}[
  >=Stealth,
  x=1cm,y=0.88cm,
  line width=0.45pt,
  line cap=round,
  line join=round,
  font=\normalsize,
  flowtext/.style={anchor=center,align=center,inner sep=0pt,
    execute at begin node={
      \setlength{\parindent}{0pt}%
      \setlength{\leftskip}{0pt}%
      \setlength{\rightskip}{0pt}%
      \everypar{}%
    }},
  every node/.style={flowtext},
  flowbox/.style={flowtext,draw,minimum height=0.68cm,inner xsep=3pt,inner ysep=2pt},
  terminal/.style={flowbox,rounded corners=0.18cm,minimum width=1.05cm},
  io/.style={flowbox,shape=trapezium,trapezium left angle=72,trapezium right angle=72,
    minimum height=0.72cm},
  process/.style={flowbox,minimum width=1.55cm},
  decision/.style={flowbox,shape=diamond,aspect=2.65,minimum width=3.95cm,minimum height=1.02cm,inner sep=0pt},
]
  \path[use as bounding box] (-3.10,0.70) rectangle (2.25,10.05);
  \node[terminal] (start) at (0,9.65) {开始};
  \node[process,minimum width=1.65cm] (init) at (0,8.35) {$a=1$};
  \node[process,minimum width=2.55cm] (update) at (0,6.95) {$a=2a+1$};
  \node[decision] (test) at (0,5.15) {$a>100?$};
  \node[io,minimum width=1.90cm] (output) at (0,3.55) {输出 $a$};
  \node[terminal] (finish) at (0,2.25) {结束};

  % Main flow and the single vertical entry into a=2a+1.
  \coordinate (loopjoin) at (0,7.55);
  \coordinate (loopleft) at (-2.65,5.15);
  \draw[->] (start.south) -- (init.north);
  \draw[->] (init.south) -- (loopjoin) -- (update.north);
  \draw[->] (update.south) -- (test.north);
  \draw[->] (test.south) -- node[flowtext,right] {是} (output.north);
  \draw[->] (output.south) -- (finish.north);

  % 否 returns left, then the horizontal arrow joins the main stem above a=2a+1.
  \draw[->] (test.west) -- node[flowtext,above] {否} (loopleft) |- (loopjoin);
\end{tikzpicture}
